\begin{textsample}{POS Dim 3 – al – Score 18.00 – al/al\_inf\_11.txt}  \label{ex:f3_pos_216}
1.
UM BREVE HISTÓRICO DA EDUCAÇÃO \textbf{INFANTIL} NO BRASIL A Educação \textbf{Infantil} tem uma história muito recente no Brasil.
Apesar das diversas indicações prescritas na legislação educacional voltadas para a educação pré-escolar nos anos 60 e 70, é apenas na década de 80 que é possível vermos o reconhecimento da educação enquanto direito da \textbf{criança} \textbf{pequena}, de zero a seis anos, sendo esta uma opção da família e dever do Estado na Constituição Federal de 1988.
Com essa conquista, fomentada por movimentos sociais, vemos fortalecer, na história da Educação \textbf{Infantil}, a busca por garantias deste direito, visando uma proposta em prática para o desenvolvimento integral do sujeito.
No decorrer do tempo, estudos e pesquisas se ampliam com o objetivo de discutir a função da \textbf{creche} e \textbf{pré-escola}, uma vez que, para a primeira, na maioria das vezes, a prática se respaldava em propostas assistencialistas, compensatórias, com ênfase nas carências, voltada para o \textbf{cuidado} da \textbf{criança} em tempo integral, filhos de mães trabalhadoras e menos favorecidas socialmente, enquanto a segunda era ofertada para uma classe social mais abastada e previa uma educação voltada para a escolarização, fortalecendo assim a cisão entre o \textbf{cuidar}.
A discussão em torno do desenvolvimento \textbf{infantil} e a importância de práticas institucionais que o favorecessem passam a ganhar destaque no meio acadêmico e a se expandir para o contexto político e prático.
Passa -se a defender, independentemente da classe social, a educação e o \textbf{cuidado} da \textbf{criança} \textbf{pequena}, de modo integrado em espaços institucionais, acessível a todas.
A elaboração e o desenvolvimento de políticas públicas numa dimensão em que cada instância federativa passaria a ocupar um lugar estratégico para que o acesso dessa faixa \textbf{etária} fosse garantido com sucesso.
Dois anos depois, em 1990, o Estatuto da Criança e do Adolescente ( ECA ) reafirma os direitos constitucionais em relação à Educação \textbf{Infantil}, sendo esse um salto que provocaria a União a definir metas para que o atendimento das \textbf{crianças} \textbf{pequenas} englobasse todos os direitos e necessidades, o que começava a apontar para um novo conceito de \textbf{infância}, \textbf{criança} e educação.
Em 1994, o Ministério da Educação e Cultura ( MEC ) publicou o documento Política Nacional de Educação \textbf{Infantil}, que estabeleceu metas como a expansão de vagas e políticas de melhoria da qualidade no atendimento às \textbf{crianças}.
Entre elas, a necessidade de qualificação dos profissionais, que resultou no documento Por uma política de formação do profissional de Educação \textbf{Infantil}.
Em 1996, com a promulgação da Emenda Constitucional que cria a Lei de Diretrizes e Bases da Educação Nacional ( LDBEN ), reconhece e regulamenta a Educação \textbf{Infantil} como primeira etapa da Educação Básica.
O artigo 62 foi pioneiro ao estabelecer a necessidade de formação para o profissional da Educação \textbf{Infantil}.
Segundo a lei, a formação do educador desse segmento deve ser “ em nível superior, admitindo -se, como formação \textbf{mínima}, a oferecida em nível médio, na modalidade Normal ”.
O texto reafirma, também, a responsabilidade constitucional dos municípios na oferta de Educação \textbf{Infantil}, \textbf{contando} com a assistência técnica e financeira da União e dos estados.
Desse modo, a Educação \textbf{Infantil} passou a ser a primeira etapa da Educação Básica, seguida pelas etapas de ensinos Fundamental e Médio.
Enquanto etapa, a \textbf{criança} de 0 a 6 anos passa a ter o direito de oferta de acesso à Educação, que a compreende numa dimensão mais ampla dentro do sistema educacional, passando a ser definida como alguém capaz de criar e estabelecer relações, um ser sócio-histórico, produtor e consumidor de cultura, inserido nela e que, portanto, não precisa apenas de \textbf{cuidado}, mas está preparado para afirmar -se enquanto sujeito de linguagem dentro das singularidades que o definem, atribuindo assim sentidos à sua experiência através de diferentes linguagens, como meio para seu desenvolvimento em diversos aspectos ( afetivos, psicológicos, cognitivos, motores e sociais ).
Em 1998, com o objetivo de oferecer parâmetros para a \textbf{manutenção} e a criação de novas \textbf{instituições} de Educação \textbf{Infantil}, o MEC publicou o documento Subsídios para credenciamento e o funcionamento das \textbf{instituições} de Educação \textbf{Infantil}.
No mesmo ano, visando a um movimento que incentivasse a elaboração de currículos de Educação \textbf{Infantil}, cuja responsabilidade foi delegada a cada \textbf{instituição} e seus profissionais, o Ministério editou o Referencial Curricular Nacional para a Educação \textbf{Infantil} ( RCNEI ), como parte dos Parâmetros Curriculares Nacionais.
Um ano depois, em 1999, o Conselho Nacional de Educação ( CNE ) publicou as Diretrizes Curriculares Nacionais para a Educação \textbf{Infantil}, constituindo -se esse documento como base legal e subsídio para elaboração e avaliação das propostas pedagógicas das \textbf{instituições} de Educação \textbf{Infantil} de todo o país.
O debate gira em torno das exigências do trabalho com \textbf{crianças} de zero a seis anos, de modo que requer professores com amplo conhecimento político, cultural e social, bem como a compreensão sobre os processos de aprendizagem e desenvolvimento humano ( BARBOSA, 2011 ), o que induz à elaboração de políticas voltadas para a melhoria da qualidade da educação, dentre estas, a formação deste profissional.
Em 2005, com o propósito de contribuir para o cumprimento da meta de Formação de Professores para Educação \textbf{Infantil}, frente à provocação definida no Plano Nacional de Educação ( PNE ), que sinalizava, através da meta 1.8, que todos os profissionais lotados como professores para trabalhar com a demanda de Educação \textbf{Infantil} tivessem formação em nível médio ( modalidade Normal ) até 2006 e 70 \% com formação em nível superior até 2011, foi criado e lançado pelo MEC o Programa de Formação Inicial para Professores em Exercício na Educação \textbf{Infantil}, conhecido como ProInfantil.
O referido curso de nível médio, a distância, na modalidade Normal, teve como público-alvo profissionais que atuavam em sala de aula da Educação \textbf{Infantil}, nas \textbf{creches} e \textbf{pré-escolas} das redes públicas — municipais e estaduais — e da rede privada, sem fins lucrativos — comunitárias, filantrópicas ou confessionais, conveniadas ou não —, sem a formação específica para o magistério.
O curso, com duração de dois anos, dispunha de material pedagógico específico para a educação a distância.
Essa foi uma tentativa que não logrou muito êxito, uma vez que a falta de continuidade, entre outros entraves, não fez dessa proposta uma ferramenta eficaz de formação.
Para a garantia da formação em nível de graduação, em maio de 2006, o Conselho Nacional de Educação ( CNE ) institui as Diretrizes Curriculares Nacionais para o Curso de Graduação em Pedagogia, licenciatura, por meio da Resolução CNE/CP nº 1, nas quais define o lócus da formação inicial para o exercício da docência na Educação \textbf{Infantil}. > A Educação \textbf{Infantil} passa a ser vista como a junção do educar e \textbf{cuidar}. \textbf{Cuidar} no sentido de que as necessidades básicas da \textbf{criança} sejam atendidas e educar porque deve oferecer à \textbf{criança} possibilidades de descobertas e aprendizados, de acordo com a Lei 9.131/95, art. 3º, que afirma : > > “ [... ] III – As \textbf{Instituições} de Educação \textbf{Infantil} devem promover em suas Propostas Pedagógicas práticas de educação e \textbf{cuidados} que possibilitem a integração entre os aspectos físicos, emocionais, afetivo-cognitivos/linguísticos e sociais da \textbf{criança}, entendendo que ela é um ser completo, total e indivisível ” ( BRASIL, 1995 ).
No que concerne às etapas da Educação Básica, vemos a reformulação do Ensino Fundamental, que passa de 8 para 9 anos, por meio da Lei nº 11.274, de 6 de fevereiro de 2006.
Isso afeta diretamente as \textbf{crianças} de 6 anos, até então atendidas pela Educação \textbf{Infantil}, que passam a ser atendidas na etapa do Ensino Fundamental.
A faixa \textbf{etária} atendida nas unidades de Educação \textbf{Infantil} passa a ser de 0 a 5 anos e 11 \textbf{meses}.
No que se refere a financiamento, a Educação \textbf{Infantil} tem um ganho substancial ao ser inserida na Lei do FUNDEB, em 2007, com ampliação do financiamento para essa etapa educacional.
A partir do princípio de regime de colaboração previsto na CF de 88, visando colaborar com os municípios em prol da qualidade na perspectiva de financiamento, a ampliação da oferta é proposta pela União a partir do Programa Nacional de Reestruturação e Aquisição de Equipamentos para a Rede Escolar Pública de Educação \textbf{Infantil} ( Proinfância ), pela Resolução nº 6, de 24 de abril de 2007, compreendendo uma das ações do Plano de Desenvolvimento da Educação ( PDE ) do MEC e depois fazendo parte do Plano de Ações Articuladas ( \textbf{PAR} ), na dimensão de infraestrutura, visando à garantia do acesso de \textbf{crianças} a \textbf{creches} e escolas, bem como à melhoria da infraestrutura física da rede de Educação \textbf{Infantil}.
Através dessa ação, a discussão em \textbf{volta} das questões e conceitos referentes a ambientes e espaços para \textbf{crianças} dessa etapa ganha padrões para suas estruturas físicas e fomenta reflexões pedagógicas para a utilização dos espaços, favorecendo o objetivo principal : a construção de uma identidade da \textbf{criança} a partir de quem ela é, com direitos de aprendizagem a serem garantidos, obrigatoriamente, pelo Estado.
Para garantir esta oferta, em novembro de 2009, por meio da Emenda Constitucional nº 59, o atendimento de \textbf{crianças} a partir de 4 anos de idade se torna obrigatório e gratuito.
Paralelo a isso, as DCNEI são atualizadas em 2009 e publicadas em 2010, mediante a necessidade de se revisar as concepções em torno da educação de \textbf{crianças} em espaços coletivos, e de seleção e fortalecimento de práticas pedagógicas mediadoras de aprendizagens e do desenvolvimento das \textbf{crianças}.
Desse modo, passam a orientar as políticas públicas e a elaboração, planejamento, execução e avaliação de propostas pedagógicas e curriculares de Educação \textbf{Infantil}.
E é nesse contexto que surgem programas suplementares ( Brasil Carinhoso, Biblioteca na Escola, Caminho da Escola ), para subsídio de materiais didáticos escolares, transporte, alimentação e assistência à saúde.
A \textbf{creche} e a \textbf{pré-escola} passaram a dispor de recursos, investimentos e políticas públicas para assegurar o acesso e permanência com algumas garantias.
Em setembro de 2012, por meio da Resolução CD/FNDE/MEC nº 19 do Programa Brasil Carinhoso, Lei 12.722/2012, que estabelece os procedimentos operacionais para a transferência obrigatória de recursos financeiros aos municípios e ao Distrito Federal, a título de apoio financeiro suplementar à \textbf{manutenção} e ao desenvolvimento da Educação \textbf{Infantil}, com eixo específico de ampliação do acesso à \textbf{creche} de \textbf{crianças} de 0 a 3, cujas famílias sejam beneficiárias do Programa Bolsa Família, em \textbf{creches} públicas ou conveniadas com o poder público.
Com esse apoio da União, principalmente as unidades de \textbf{creche} passaram a ser reorganizadas, vislumbrando a erradicação da miséria, assistindo as \textbf{crianças} em condições de vulnerabilidade social.
Preconiza -se que, por meio desta ação de transferência de recursos suplementares, a oferta se ampliasse e houvesse uma melhora nas condições e na qualidade da educação e a redução da desigualdade educacional para as \textbf{crianças} de 0 a 3 anos.
Diante deste contexto de avanços legais e políticas públicas, é possível identificar a mudança da compreensão das especificidades em torno da \textbf{criança} e da \textbf{infância}, conforme defendido nas teorias do desenvolvimento, evidenciando a complexidade do desenvolvimento \textbf{infantil}, tal como descreve Vigotsky ( 1989 ) : > Acreditamos que o desenvolvimento da \textbf{criança} é um processo dialético complexo caracterizado pela periodicidade, desigualdade no desenvolvimento de diferentes funções, metamorfose ou transformação qualitativa de uma forma em outra, imbricamento de fatores internos e externos, e processos adaptativos que superam os impedimentos que a \textbf{criança} encontra ( p. 51 ).
Percebe -se a complexidade dessa etapa e de sua importância para o desenvolvimento do sujeito.
Iniciou -se uma busca de implementação de novas políticas públicas no Brasil que garantisse esse desenvolvimento integral, como também o alinhamento entre programas, projetos e planos federais, estaduais e municipais com metas, estratégias e propostas pedagógicas que favorecessem desde a infraestrutura à formação do professor, a uma Educação \textbf{Infantil} que proporcionasse, mesmo em meio à diversidade que compreende nosso país, os avanços que oportunizam saberes e habilidades que se constituem direitos de aprendizagem nessa etapa da vida.
No ano de 2017, o MEC oferta uma proposta de formação para professores de \textbf{pré-escola}, incluindo no PNAIC a demanda de 4 a 5 anos que compreende a etapa de Educação \textbf{Infantil}.
Com essa ação, visualiza -se a importância dessa etapa para o processo de alfabetização, contudo, tendo claro que a proposta não se caracteriza como um retrocesso para a escolarização nessa etapa e sim, frente à proposta de desenvolvimento integral da \textbf{criança} de 4 e 5 anos, desenvolver no docente as competências que o farão hábeis para mediar a aquisição de habilidades, através das vivências e interações a serem incluídas nos planos de aula e consequentemente desfrutadas no fazer pedagógico pelas \textbf{crianças} \textbf{pequenas}.
Dentro desse movimento, a Base Nacional Comum Curricular ( BNCC ), prevista na Constituição de 1988, na LDB de 1996 e no Plano Nacional de Educação de 2014, foi homologada em 22 de dezembro de 2017 pelo Conselho Nacional de Educação ( CNE ), constando como resolução para nortear o currículo de Educação \textbf{Infantil} em todo o país, tendo por eixos estruturantes as interações e \textbf{brincadeiras}, “ experiências nas quais as \textbf{crianças} podem construir e apropriar -se de conhecimentos por meio de suas ações e interações com seus \textbf{pares} e com os \textbf{adultos}, o que possibilita aprendizagens, desenvolvimento e socialização ” ( BRASIL, 2017, p. 37 ).
Para tanto, o documento impõe a necessidade de intencionalidade educativa desde a \textbf{creche} até as turmas de 5 anos e 11 \textbf{meses}, que garanta direitos de aprendizagens específicos dessa faixa \textbf{etária}.
Esta intencionalidade consiste : > “ [... ] na organização e proposição, pelo educador, de experiências que permitam às \textbf{crianças} conhecer a si e ao outro e de conhecer e compreender as relações com a natureza, com a cultura e com a produção científica, que se traduzem nas práticas de \textbf{cuidados} pessoais ( alimentar -se, vestir -se, higienizar -se ), nas \textbf{brincadeiras}, nas experimentações com materiais variados, na aproximação com a literatura e no encontro com as pessoas ” ( BRASIL, 2017, p. 37 ).
Em 2018, a educação brasileira vive um momento de implementação da Base, conduzindo cada unidade federativa à revisão ou construção de seus Referenciais Curriculares Estaduais, à luz do referido documento normativo, que subsidiarão os Currículos dos municípios e consequentemente os Projetos Políticos Pedagógicos ( PPPs ) das \textbf{creches} e \textbf{pré-escolas} das redes públicas municipais e estaduais — e rede privada, sem fins lucrativos — comunitárias, filantrópicas ou confessionais — conveniadas ou não.
Portanto, esse é um momento de adequação dos currículos à proposta educacional vigente, compreendendo 60 \% ao que está contemplado na Base e 40 \% ao que concerne à singularidade da oferta de cada região, localidade e especificidade existente em nosso país.
Com esse pano de fundo, o estado de Alagoas, através de uma equipe formada com representações da UNDIME e SEDUC, assume a responsabilidade de revisar o Referencial Curricular da Educação Básica da Rede Estadual de Ensino de Alagoas — Educação \textbf{Infantil} — tendo por proposta ampliar as discussões, atualizando nomenclaturas, conceitos e preceitos a partir da BNCC, vislumbrando referenciar os sistemas e redes de nosso Estado rumo à implantação e implementação dos campos de experiências e direitos de aprendizagens, que comporão os Currículos a partir de então, em busca de garantir uma gestão escolar e de sala de aula que respeite o direito da \textbf{criança} de ser \textbf{criança}, oportunizando nesse espaço a socialização de saberes e fazeres pedagógicos que proporcionarão aos profissionais que atendem a etapa em questão ( Técnicos de Educação \textbf{Infantil} das Secretarias Municipais de Educação, Coordenadores pedagógicos, professores, auxiliares de sala, entre outros profissionais envolvidos ) um documento para subsídio de práticas significativas e singulares em respeito à complexidade que envolve esse momento do desenvolvimento da \textbf{criança}.

% matched lemmas: adulto, brincadeira, contar, creche, criança, cuidado, cuidar, etário, infantil, infância, instituição, manutenção, mês, par, pequeno, pré-escola, volta
\end{textsample}
